% ============================================================================
% 10f-appendix-supplychain.tex -- Supply-chain gate and physical sandbox
% Evidence: acceptance/supply-chain/2026-10-28-供应链防御验收.md / .json
%   policy:   contracts/supply-chain-policy.json
%             (10 capabilities, 10 scan rules, 19 error codes, 11-step precedence)
%   impl:     plugin/dsh-lingshu/supply-chain.js      (pure decision functions)
%             plugin/dsh-lingshu/sandbox-host.mjs    (parent / RPC / kill)
%             plugin/dsh-lingshu/sandbox-child.mjs   (Node permission model)
%             engine/supply_chain.py                 (Python reference, parity checked)
%   measured: sim 59/59 assertions; integration 10/10; python test 172 assertions;
%             py<->js parity 32 cases, 0 field differences (2026-10-28, Node 24.21.0)
% ============================================================================

\section{Supply-Chain Gate and Physical Sandbox}
\label{sec:app:supplychain}

\paragraph{Why an infrastructure layer.}
The control plane analysed in Section~\ref{sec:actuator} and
Appendix~\ref{app:tuca} governs calls that an agent \emph{decides} to make. It
cannot contain a plugin whose module body executes with host privileges the
moment it is imported: by the time any tool rule fires, a malicious
\texttt{import} has already run. We therefore separate two threats that are
often conflated: (i) a third-party cognition plugin that is malicious or
compromised --- supply-chain poisoning; and (ii) a model attempting an
out-of-policy tool call --- delegated to the T-UCA gate. This appendix covers
(i) only. The two layers meet at the audit chain and at one shared capability
vocabulary; they deliberately do not share a verdict vocabulary, because a
machine cannot ask a human a question during \texttt{import}.

\paragraph{Threat model and the fail-closed rule.}
Two facts from our own host audit motivate the design: the host does not
verify a plugin's declared capabilities (hypothesis H4 was falsified in
\texttt{research/02 \S4}: the patch layer can rewrite any config), and a
plugin's on-disk package can be modified by anything that can write the
profile directory. The gate therefore makes no assumption of good behaviour.
Three checks run in order: manifest validation, lexical source scan, and
process-level sandboxing. The verdict is binary: \texttt{allow\_sandboxed} or
\texttt{deny}. A deny means the module is never imported, no hook is
registered, one \texttt{supply\_chain} audit event is written, and the loader
is suspended (later plugins are skipped). There is no \texttt{degrade} in this
layer by design: degradation is a human-in-the-loop decision and belongs to
the control plane (ADR-017); infrastructure must not trade safety for
availability.

\paragraph{Manifest contract.}
A plugin ships \texttt{plugin.manifest.json} with: \texttt{manifest\_version}
(\texttt{lingshu.plugin-manifest/v1}); \texttt{id} and \texttt{version};
\texttt{entry} (relative path, parent traversal rejected); \texttt{exports}
(callable names, at most 32); \texttt{capabilities.required} /
\texttt{capabilities.optional}; \texttt{fs.read} / \texttt{fs.write} scopes
(relative, must stay inside the plugin directory); and
\texttt{integrity.entry\_sha256}. The capability vocabulary is inherited
verbatim from the \texttt{dsh-trust-check} ecosystem (ten names, 3841 scanned
plugins). Only three are grantable: \texttt{fs-read}, \texttt{fs-write},
\texttt{llm}. The other seven --- \texttt{network}, \texttt{shell},
\texttt{credentials}, \texttt{dynamic-code}, \texttt{subagent},
\texttt{env}, \texttt{host-runtime} --- are \emph{never} grantable; declaring
one is itself a denial. A minimal manifest:

{\footnotesize\begin{verbatim}
{ "manifest_version": "lingshu.plugin-manifest/v1",
  "id": "demo-emotion", "version": "0.1.0",
  "entry": "index.mjs", "exports": ["ping"],
  "capabilities": { "required": ["fs-read"] },
  "fs": { "read": ["data"] },
  "integrity": { "entry_sha256": "<64 hex>" } }
\end{verbatim}}

\paragraph{Lexical pre-filter, not the boundary.}
Ten rules (SC-01\ldots SC-10) map API families to capabilities:
write-side \texttt{fs} calls, read-side \texttt{fs} calls,
\texttt{child\_process}, network modules and \texttt{fetch}, credential
material, dynamic code (\texttt{eval}, \texttt{Function}, \texttt{vm},
computed \texttt{import}), worker threads, native addons and
\texttt{process.binding}, \texttt{process.env}, and obfuscation signals
(base64/hex/atob). Line and block comments are stripped before matching, so
prose that names an API does not trip a rule; string literals are
deliberately \emph{not} stripped, which is conservative (a documentation
string containing \texttt{fs.writeFileSync} is denied rather than trusted).
A finding whose capability is not declared yields
\texttt{undeclared\_usage} and a deny. We make no completeness claim: names
built by concatenation, computed dynamic imports, and
template-string interpolation evade lexical matching. That is exactly why the
scan runs \emph{before} the sandbox rather than instead of it.

\paragraph{Verdict precedence and error codes.}
Blockers are ordered and the first one is reported as the reason:
\texttt{sandbox\_unavailable}, \texttt{manifest\_missing},
\texttt{manifest\_unreadable}, \texttt{manifest\_invalid},
\texttt{entry\_missing}, \texttt{entry\_too\_large},
\texttt{integrity\_mismatch}, \texttt{fs\_scope\_escape},
\texttt{unknown\_capability}, \texttt{forbidden\_capability},
\texttt{undeclared\_usage} (SC-E01\ldots SC-E11); runtime failures continue
the same namespace with SC-E12\ldots SC-E19 (start failed, ready timeout,
call timeout, crash, spawn failed, handshake failed, method not allowed,
plugin error) and SC-E20 (skipped after suspension). Two codes are
intentionally shareable rather than distinct: \texttt{unknown\_capability}
is normally shadowed by \texttt{manifest\_invalid}, and
\texttt{fs\_scope\_escape} is reachable through platform-specific
traversal (e.g.\ backslash segments that the manifest validator does not
split on). The whole policy is one file,
\texttt{contracts/supply-chain-policy.json}, loaded at startup by both the
JavaScript implementation and the Python reference; a missing or corrupt
policy degrades to a deny-all fallback rather than to permissive defaults.

\paragraph{Physical sandbox.}
An allowed plugin is executed in a child process under Node's permission
model. Flags are derived mechanically from the manifest: \texttt{--permission}
always; \texttt{--allow-fs-read} for the plugin directory (module loading
requires it) plus each declared read scope; \texttt{--allow-fs-write} only for
declared write scopes; \texttt{--no-addons},
\texttt{--disable-proto=throw},
\texttt{--disallow-code-generation-from-strings} and
\texttt{--max-old-space-size=128} unconditionally. No
\texttt{--allow-child-process} and no \texttt{--allow-worker} flag is ever
emitted. The parent and child speak newline-delimited JSON over a loopback TCP
socket with a one-shot random token; only the exports declared in the manifest
are callable; ready timeout is 5000\,ms and per-call timeout 2000\,ms; a
timeout or uncaught violation kills the child and is reported as a deny.
The transport is TCP rather than stdio pipes for an observed reason: under the
host's own file sandbox, spawning with piped stdio fails with
\texttt{EPERM}, while \texttt{stdio: inherit/ignore} works --- a constraint we
record rather than hide.

\paragraph{Measured enforcement (2026-10-28, Node 24.21.0).}
The plugin-side suite (\texttt{supply-chain-sim.mjs}) passes 59/59
assertions against real child processes, including: an out-of-scope file
write and an out-of-scope read both fail with \texttt{ERR\_ACCESS\_DENIED};
\texttt{child\_process} is denied by the permission model; a write performed
through a concatenated property name (\texttt{fs["wri"+"teFileSync"]}) is
\emph{not} flagged by the lexical scan (zero SC-01 findings) and is still
denied at runtime, with the target file absent afterwards --- the concrete
case where the sandbox catches what the scanner misses; an import-time escape
fails the load with SC-E12 and leaves no file; a 300\,ms call timeout kills
the child and returns SC-E14; a declared \texttt{network} capability is denied
before any process is spawned. The index-level integration harness applies the
real plugin with two third-party directories: both manifests are parsed, the
benign plugin runs in a child whose pid differs from the host, the malicious
plugin is denied and its marker file is never created, the loader is
suspended, and the read-only status tool reports the suspended state (10/10
assertions). Cross-language parity covers 32 decision cases (all error paths,
comment stripping, scope escapes, optionals): Python and JavaScript agree on
every contract field, 0 differences. The Python acceptance file totals 172
assertions. Evidence, including raw JSON, is in
\texttt{acceptance/supply-chain/}.

\paragraph{Stated limits.}
(1) The permission model does not cover the network: a \texttt{connect} under
\texttt{--permission} succeeds (verified). Network egress is therefore
prevented only by never granting \texttt{network} and by the lexical
pre-filter; OS-level egress control is future work and is not claimed.
(2) The plugin directory must remain readable for module loading, so a plugin
that requests no \texttt{fs-read} can still read files inside its own
directory; declared scopes add to, they do not restrict, that baseline.
(3) Integrity is a pinned SHA-256: it detects drift between audit and load, not
a malicious author who ships a matching hash. No signature/PKI is implemented.
(4) Scan blind spots are listed above; the sandbox is the backstop, and the
sandbox itself is one process-level boundary, not a VM. (5) WASM was
considered and not adopted: existing plugins are ESM packages, and a
subprocess already isolates the vulnerable native surface at a small migration
cost; a WASM runtime is a future extension, not a current claim. (6) The gate
has been exercised against synthetic plugins generated by the test harness,
not yet against plugins from a live registry; the acceptance record says so
explicitly.

\paragraph{Relationship to the control plane.}
The two layers close different gaps and share obligations. The infrastructure
layer contributes: capability minimisation expressed in the same ten-name
vocabulary as the host scanner; a binary, pre-execution decision with an audit
record; and a process boundary that survives a compromised model. The control
plane (Section~\ref{sec:actuator}, Appendix~\ref{app:tuca}, ADR-017)
contributes the explainable escalation path when a high-risk \emph{tool} call
is degraded rather than blocked. Neither layer claims the other's job; a
plugin that passes the gate is still subject to every T-UCA rule for the tools
it calls, and a model that obeys every tool rule can still be stopped from
importing hostile code. That separation --- not any single check --- is what
``defence in depth'' means in this system.
